% ==============================================================================
% Capítulo de Conclusão
% ==============================================================================

\chapter{CONCLUSÃO}\label{cap:conclusao}

Este trabalho investigou como as despesas dos municípios brasileiros respondem ao calendário eleitoral e ao alinhamento político do prefeito com o governador e com o presidente. Foram usados dados do Siconfi, do TSE e do IBGE para cerca de 5,6 mil municípios entre 2013 e 2025, em dois modelos de efeitos fixos: o primeiro mede o ciclo eleitoral, com erros-padrão de Driscoll e Kraay; o segundo mede o efeito do alinhamento político, com efeitos fixos de município e de ano. A despesa entra em logaritmo, de modo que os efeitos se leem como variações percentuais, e a especificação em nível foi mantida como teste de robustez e para a comparação com \textcite{sakurai2009}.

O principal resultado é que o ciclo eleitoral aparece sobretudo na composição do gasto. No ano eleitoral, a despesa em Habitação e Urbanismo é cerca de 27\% maior, a de Transporte, cerca de 19,5\% maior, e a de Saúde e Saneamento, cerca de 4,2\% maior. Esses três aumentos ocorrem nas eleições de 2016 e de 2024, nas duas formas funcionais e com os dois tipos de erro-padrão, e são os resultados mais sólidos do trabalho. Duas ressalvas: o efeito em Saúde e Saneamento deixa de ser significativo quando o PIB nacional é trocado pelo municipal (R7), e a magnitude em Transporte deve ser lida com cautela, pelos valores pequenos de parte dos municípios. A participação das áreas na despesa total também muda: no ano eleitoral, Habitação e Urbanismo e Transporte ganham peso, enquanto Administração e Educação e Cultura perdem. Esse padrão é compatível com um deslocamento de recursos para os gastos mais visíveis ao eleitor, como obras urbanas e transporte.

Na despesa total, o ano eleitoral está associado a um gasto cerca de 5,3\% maior, mas esse resultado pede cautela. O efeito é de cerca de 3,3\% em 2016 e de 6,7\% em 2024 e, quando a amostra termina em 2023, deixa de ser significativo na especificação principal. Em 2024, quase todas as áreas sobem ao mesmo tempo, inclusive a Administração, um gasto pouco visível. Esse aumento generalizado admite duas leituras: é compatível com o forte crescimento das receitas municipais no período e também com uma estratégia eleitoral, já que os investimentos municipais cresceram fortemente em 2024, alta que o Ipea associa ao ano eleitoral \parencite{santana2025}. O desenho empírico, no entanto, não permite separar as duas explicações. Outros resultados também devem ser lidos com cuidado: o efeito em Educação e Cultura depende do tratamento dado a 2020, ano de eleição e de pandemia; o de Assistência Social aparece só em 2024; e o de Agricultura depende da forma funcional.

Quanto às coalizões, prefeitos da coligação do governador gastam cerca de 0,37\% a mais no total, um efeito pequeno, mas que se mantém em todos os testes que incluem esse coeficiente, e cerca de 0,7\% a mais em Administração. Como as transferências correntes já estão entre os controles, esse efeito pode refletir canais de apoio estadual não captados pelas transferências correntes, como as transferências de capital. O alinhamento com o presidente, por sua vez, não tem efeito sobre a despesa total na amostra completa, e o resultado é sensível ao período: na amostra de 2013 a 2023, o coeficiente fica negativo, de cerca de $-$0,3\%. A comparação entre os dois modelos mostra a importância do efeito fixo de ano: como a proporção de prefeitos alinhados ao presidente muda de forma brusca entre governos, sem esse controle o alinhamento se confunde com o período.

Os resultados são compatíveis com os modelos de sinalização de \textcite{rogoff1990}, em que o governante expande, antes da eleição, os gastos de rápida percepção pela população, e com o ciclo de composição descrito por \textcite{drazen2010}. Os dados, contudo, não permitem distinguir se esse comportamento decorre de oportunismo ou de sinalização de competência, já que as duas abordagens preveem aumento dos gastos visíveis antes da eleição. Os resultados também confirmam parte dos achados de \textcite{sakurai2009} para 1990--2005, que encontrou aumento no ano eleitoral em Saúde e Saneamento, Habitação e Urbanismo e Transporte, o que indica que esse padrão persiste no período recente. O efeito do alinhamento com o governador é coerente com \textcite{ferreira2007}, que mostram que municípios governados por aliados recebem mais transferências voluntárias.

O trabalho tem limitações. O período cobre 13 anos, e o efeito do ano eleitoral é identificado por apenas duas eleições, porque 2020 coincide com a pandemia; com poucos períodos, os erros-padrão de Driscoll e Kraay perdem precisão. A despesa é medida no estágio do pagamento, de modo que não capta as despesas empenhadas e não pagas no exercício, inscritas em restos a pagar, que \textcite{vicente2012} apontam como instrumento de uso estratégico no ano eleitoral. A classificação funcional é informada pelo próprio município, e funções como Transporte e Agricultura têm registro irregular. Em Comunicações, os valores são tão pequenos que o coeficiente não se lê como variação percentual, o que impediu a interpretação desse grupo. Além disso, as transferências correntes incluem mais do que os repasses da União e dos estados, a ideologia do partido do prefeito não foi considerada, e o PIB de 2024 e 2025 ainda é preliminar. A série oficial de população também muda de base em 2022, o que afeta os valores per capita de 2022 e 2023; o teste R17, com esses dois anos interpolados, indica que isso não altera os resultados principais.

Essas limitações indicam caminhos para trabalhos futuros. O uso dos estágios de empenho e de inscrição em restos a pagar permitiria verificar se parte do ciclo ocorre fora do pagamento. A classificação por natureza da despesa permitiria separar investimentos de custeio. Dados de transferências voluntárias e de capital permitiriam testar o canal do alinhamento com o governador. A comparação entre prefeitos em primeiro e em segundo mandato, ou entre municípios com diferentes níveis de competição eleitoral, permitiria avaliar se o ciclo depende dos incentivos à reeleição. Por fim, a inclusão das próximas eleições municipais ajudaria a esclarecer se o efeito sobre a despesa total observado em 2024 é recorrente ou específico daquele período.
